\centerline{\bf \Huge Лемма 255}
\centerline{\bf \Huge }

\begin{flushright}
        \scriptsize{А лемма номер 255, на телогреечке печать2}
\end{flushright}

\problem{1}
Пусть $G_b$ и $G_c$-точки касания вписанной окружности со сторонами $A C$ и $A B$, а точки $M_a$ и $M_b$ - середины $B C$ и $A C$. Докажите, что прямые $G_b G_c, B I, M_a M_b$ пересекаются в точке на окружности, построенной на $B C$ как на диаметре.

\problem{2}
$\bigl[\textbf{ВсОШ, регион, 9.2, 2013}\bigr]$
 Окружность, вписанная в прямоугольный треугольник $A B C$ с гипотенузой $A B$, касается сторон $B C, C A, A B$ в точках $A_1, B_1, C_1$ соответственно. Пусть $B_1 H$-высота треугольника $A_1 B_1 C_1$. Докажите, что точка $H$ лежит на биссектрисе угла $C A B$.

\problem{3}
$\bigl[\textbf{Олимпиада СПб, 1999}\bigr]$
 В неравнобедренном треугольнике $A B C$ проведены биссектрисы $A A_1$ и $C C_1$, кроме того, отмечены середины $K$ и $L$ сторон $A B$ и $B C$ соответственно. На прямую $C C_1$ опущен перпендикуляр $A P$, а на прямую $A A_1-$ перпендикуляр $C Q$. Докажите, что прямые $K P$ и $L Q$ пересекаются на стороне $A C$.

\problem{4}
$\bigl[\textbf{MМО, 1999}\bigr]$
Вписанная окружность треугольника $A B C$ касается сторон $A B$ и $B C(A B>B C)$ в точках $N$ и $M$ соответственно. $P Q$ - средняя линия треугольника $A B C$, параллельная $A B, K$-точка пересечения $M N$ и $P Q$. Докажите, что $A K-$ биссектриса угла $B A C$.

\problem{5}
В треугольнике $A B C$ выполняется равенство $3 A C=A B+B C$. Вписанная в треугольник окружность касается сторон $A B$ и $B C$ в точках $K$ и $L$ соответственно; $D K$ и $E L$ - ее диаметры. Докажите, что точки пересечения прямых $A E$ и $C D$ с прямой $K L$ равноудалены от середины отрезка $A C$.

\problem{6}
Дан треугольник $A B C$. На продолжении стороны $B C$ за точку $C$ отмечена точка $X$. Окружности, вписанные в треугольники $A B X$ и $A C X$, пересекаются в точках $P$ и $Q$. Докажите, что все прямые $P Q$ проходят через одну и ту же точку, не зависящую от выбора точки $X$.